% ============================================================
% P11 — Induktiver Hilbert-Grenzraum mit quellenkanonischer
%        arithmetischer Metrik
% ============================================================
% STATUS: PASS-A ACTIVE — kein SYN, kein Seal
% Letzter stabiler Knoten: C5e (2026-08-10)
% Offener Kanal: Ungerader relativer Terminal-Transport (C6ff.)
% Dieses Manuskript ist der kanonische Sammelpunkt der
% derzeit gültigen Architektur. Es enthält noch keine
% vollständigen Beweise — diese liegen in audits/.
% ============================================================

\documentclass[12pt,a4paper]{amsart}

% --- Pakete ---
\usepackage[utf8]{inputenc}
\usepackage[T1]{fontenc}
\usepackage{amsmath,amssymb,amsthm}
\usepackage{mathtools}
\usepackage{hyperref}
\usepackage{cleveref}
\usepackage[margin=2.5cm]{geometry}

% --- Theorem-Umgebungen ---
\newtheorem{theorem}{Theorem}[section]
\newtheorem{proposition}[theorem]{Proposition}
\newtheorem{lemma}[theorem]{Lemma}
\newtheorem{corollary}[theorem]{Corollary}
\newtheorem{definition}[theorem]{Definition}
\newtheorem{remark}[theorem]{Remark}
\newtheorem{construction}[theorem]{Construction}

% --- Operatoren und Notation ---
\newcommand{\KXR}{K_{X,R}}
\newcommand{\KXRp}{K_{X,R}^{+}}
\newcommand{\KXRm}{K_{X,R}^{-}}
\newcommand{\JRS}{J_{R,S}}
\newcommand{\JRSp}{J_{R,S}^{+}}
\newcommand{\GRT}{G_{R,T}}
\newcommand{\GRTp}{G_{R,T}^{+}}
\newcommand{\GammaRp}{\Gamma_R^{+}}
\newcommand{\qRT}{q_{R,T}}
\newcommand{\qRTp}{q_{R,T}^{+}}
\newcommand{\qGammaR}{q_{\Gamma,R}^{+}}
\newcommand{\WRST}[1]{W_{R,S,#1}^{[T]}}
\newcommand{\WRSGamma}[1]{W_{R,S,#1}^{[\Gamma]}}
\newcommand{\betaR}{\beta_R}
\newcommand{\BR}{B_R}
\newcommand{\CRG}{C_{\Gamma,R}}
\newcommand{\HR}{H_R}
\newcommand{\LR}{L_R}
\newcommand{\Mosco}{\xrightarrow{\mathrm{M}}}
\newcommand{\sconv}{\xrightarrow{\mathrm{s}}}

% ============================================================
\begin{document}
% ============================================================

\title{Induktiver Hilbert-Grenzraum mit quellenkanonischer
  arithmetischer Metrik}

\author{Sebastian Schmalnauer}

\address{Objekt-X-Programm, P11}

\email{sebastian.schmalnauer@gmx.at}

\date{Stand: 10. August 2026 (PASS-A ACTIVE)}

\keywords{Hilbert-Grenzraum, induktives System, arithmetische Metrik,
  Boundary-Jet, Feshbach-Kolligation, Mosco-Konvergenz,
  gerader Terminaltransport, Riemannsche Hypothese}

\subjclass[2020]{11M26, 46C05, 47A10, 46E22}

\begin{abstract}
Wir konstruieren einen induktiven Hilbert-Grenzraum $K_X$ mit
quellenkanonischer arithmetischer Metrik als spektralen Träger
für die Weil-Explizitformel. Das induktive System
$\{\KXR, \JRS\}$ wird mit einem kohärenten Gram-Kernel
$q_R^X(f) = \|\CRG^{1/2} f\|^2 + \|\LR \HR^* f\|^2$
ausgestattet. Die quellenkanonische Metrik ist flach im endlichen
Bereich (C2). Eine Paritätszerlegung trennt den geraden Sektor
(vollständig kontrolliert: Mosco-Konvergenz, starke
Resolventen-Konvergenz, kohärenter isometrischer Terminaltransport)
vom ungeraden Sektor (offen: relativer Terminal-Transport via
Boundary-Jet). Der vorliegende Text dokumentiert den derzeit
gültigen Architekturstand; vollständige Beweise liegen im
\texttt{audits/}-Verzeichnis des Repositoriums.
\end{abstract}

\maketitle

\tableofcontents

% ============================================================
\section{Introduction}
\label{sec:intro}
% ============================================================

% TODO: Motivation — quellenerstes Kopplungsproblem.
% Warum reicht L^2(R) nicht (B2-B, P03)?
% Warum kein translationsinvarianter Operator (C1y)?
% Warum kein S_p für p < infty (B2-A)?
% Ziel: spektraler Träger für Weil-Explizitformel.

\begin{remark}[Status]
  Dieser Abschnitt ist noch nicht ausgearbeitet.
  Abhängigkeiten: P03 (Haar-L$^2$-Ausschluss), B2-A, B2-B, C1y.
\end{remark}

% ============================================================
\section{Finite source graph spaces}
\label{sec:source-graphs}
% ============================================================

% TODO: Definition der endlichen Quellgraphen K_{X,R}.
% Gram-Realisierung q_R^X.
% Bounded injective transitions J_{R,S}.
% Kokyklus-Bedingung.
% Hauptquelle: C1 (Gramkernel-Konstruktion).

\begin{definition}[Kohärentes induktives System]
  \label{def:ind-system}
  % Placeholder — aus C1 zu befüllen.
  Es sei $\{\KXR, \JRS\}_{R < S}$ ein induktives System von
  Hilberträumen mit beschränkten injektiven Transitionen und
  Gram-Realisierung
  \[
    q_R^X(f) = \|\CRG^{1/2} f\|^2 + \|\LR \HR^* f\|^2.
  \]
\end{definition}

\begin{remark}[Provenienz]
  Konstruktion in Audit-Knoten~C1; Gram-Kernel-Vollständigkeit
  in~C1b--C1z.
\end{remark}

% ============================================================
\section{BC-valued incidence and prime-power geometry}
\label{sec:bc-geometry}
% ============================================================

% TODO: BC-GCD-Geometrie, prime-power Labelstruktur.
% Hauptquelle: C1k, C1k2.

\begin{remark}[Status]
  Abschnitt offen. Abhängigkeiten: C1k, C1k2, P09.
\end{remark}

% ============================================================
\section{Feshbach colligation and future metrics}
\label{sec:feshbach}
% ============================================================

% TODO: Feshbach-/Graphraumapparat.
% Zukunftsmetriken: absolute Divergenz im ungeraden Sektor.
% Hauptquelle: C1zB1, C1zB2, C3.

\begin{remark}[Status]
  Abschnitt offen. Abhängigkeiten: C1zB1, C1zB2, C1zB2C3.
\end{remark}

% ============================================================
\section{Finite terminal-gauge trivialization}
\label{sec:terminal-gauge-finite}
% ============================================================

% TODO: C2 — Metrik-Kokyklus ist flach.
% W_{R,S}^{[T]} isometrisch für jedes endliche T.
% Das ist die finite Terminal-Gauge-Struktur.

\begin{theorem}[Flache finite Terminal-Gauge, C2]
  \label{thm:flat-gauge}
  % Placeholder.
  Für jedes endliche $T$ ist der Metrik-Kokyklus
  $W_{R,S}^{[T]}$ isometrisch.
\end{theorem}

\begin{remark}[Provenienz]
  Beweis in Audit-Knoten~C2.
\end{remark}

% ============================================================
\section{Boundary asymptotics and parity decomposition}
\label{sec:parity}
% ============================================================

% TODO: Kanonischer Boundary-Jet (beta_R^{(0)}, beta_R^{(1)}, ...).
% Exhaustion des ungeraden Source-Sektors:
%   \bigcap_m ker(beta_R^{(m)}) = K_{X,R}^+.
% Hauptquelle: C3, C4, C5.

\begin{theorem}[Boundary-Jet exhaustiert ungeraden Sektor, C4]
  \label{thm:boundary-jet}
  % Placeholder.
  Es gilt
  \[
    \bigcap_{m \geq 0} \ker \betaR^{(m)} = \KXRp.
  \]
\end{theorem}

\begin{remark}[Provenienz]
  Beweis in Audit-Knoten~C4 (kein endlichdimensionaler Randquotient).
\end{remark}

% ============================================================
\section{Even sector: Mosco convergence of quadratic forms}
\label{sec:even-mosco}
% ============================================================

% TODO: sigma_T(J_{R,T} f) = O_f(T^{-1}) -> 0 für glattes gerades f.
% Mosco-Konvergenz q_{R,T}^+ ->^M q_{Gamma,R}^+.
% Starke Resolventenkonvergenz G_{R,T}^+ -> Gamma_R^+.
% Hauptquelle: C5c, C5d.

\begin{theorem}[Mosco-Konvergenz gerader Sektor, C5d]
  \label{thm:mosco-even}
  % Placeholder.
  Es gilt
  \[
    \qRTp \Mosco \qGammaR
    \quad \text{und} \quad
    \GRTp \sconv \GammaRp.
  \]
\end{theorem}

\begin{remark}[Provenienz]
  Beweis in Audit-Knoten~C5c--C5d.
\end{remark}

% ============================================================
\section{Even sector: coherent terminal-gauge limit}
\label{sec:even-terminal}
% ============================================================

% TODO: Hauptresultat C5e.
% W_{R,S,+}^{[T]} ->^s W_{R,S,+}^{[Gamma]} stark.
% W_{R,S,+}^{[Gamma]} = (Gamma_S^+)^{1/2} J_{R,S}^+ (Gamma_R^+)^{-1/2}.
% Isometrisch + kokyklisch: W_{S,U,+}^{[Gamma]} W_{R,S,+}^{[Gamma]} = W_{R,U,+}^{[Gamma]}.
% Trick: Moving-Test-Vektor, keine globale Quadratwurzel-Konvergenz nötig.

\begin{theorem}[Gerader Terminal-Gauge-Limes, C5e]
  \label{thm:even-terminal}
  Für jedes feste $R < S$ gilt auf dem gesamten geraden Graphraum
  $\KXRp$:
  \[
    \WRST{+} \sconv \WRSGamma{+},
  \]
  wobei
  \[
    \WRSGamma{+} = (\Gamma_S^+)^{1/2} \JRSp (\GammaRp)^{-1/2}.
  \]
  Der Limes-Operator $\WRSGamma{+}$ ist isometrisch und
  erfüllt die Kokyklus-Bedingung
  \[
    W_{S,U,+}^{[\Gamma]} \, W_{R,S,+}^{[\Gamma]} = W_{R,U,+}^{[\Gamma]}.
  \]
\end{theorem}

\begin{remark}[Firewall]
  Die globale Konvergenz $(G_{S,T}^+)^{1/2} \to (\Gamma_S^+)^{1/2}$
  bleibt unbewiesen und wird nicht benötigt.
  Der Beweis verwendet einen Moving-Test-Vektor
  $y_T = B_T^{-1/2} g$, der die Target-Quadratwurzel eliminiert.
\end{remark}

\begin{remark}[Provenienz]
  Vollständiger Beweis (816 Zeilen) in
  \texttt{audits/AUDIT-2026-08-10\_P11\_C1zB2C5e\_EvenTerminalGauge\_TargetSquareRootConvergence.md},
  Commit \texttt{f65bf0b}.
\end{remark}

% ============================================================
\section{Odd sector: boundary jet and relative transport}
\label{sec:odd-open}
% ============================================================

% TODO: OFFENER KNOTEN.
% Absolute Zukunftsmetrik divergiert (gesiegelt).
% Boundary-Jet B_R(z; f) ist das kanonische analytische Objekt.
% W_{R,S,-}^{[T]} Cauchy-Konvergenz: offen.
% Ungerader relativer Terminal-Transport: nächster Hauptknoten (C6ff.).

\begin{remark}[Offener Status]
  Dieser Abschnitt entspricht dem mathematisch offenen Knoten
  des Programms (Stand 10. August 2026). Die absolute
  Zukunftsmetrik im ungeraden Sektor divergiert (Audit-Knoten
  C1zB2C, B2-B). Der Boundary-Jet $\BR(z; f)$ ist das
  kanonische analytische Objekt für den relativen Transport;
  die Cauchy-Konvergenz von $W_{R,S,-}^{[T]}$ ist offen.
\end{remark}

% ============================================================
\section{Object-X implications and remaining obstruction}
\label{sec:obstruction}
% ============================================================

% TODO: Verbindung zur Weil-Explizitformel.
% Verbindung zur Riemannschen Hypothese.
% Was fehlt noch: ungerader Kanal + spektrale Einbettung.

\begin{remark}[Status]
  Abschnitt offen. Hängt von Abschnitt~\ref{sec:odd-open} ab.
\end{remark}

% ============================================================
% Literatur (Placeholder — aus REFERENCES.md zu befüllen)
% ============================================================

\begin{thebibliography}{99}

\bibitem{P01} Schmalnauer, S. \textit{[P01 — Titel aus papers/]},
  Objekt-X-Programm, 2025.

\bibitem{P03} Schmalnauer, S. \textit{[P03 — Haar-L$^2$-Ausschluss]},
  Objekt-X-Programm, 2025.

\bibitem{P05} Schmalnauer, S. \textit{[P05 — SYN]},
  Objekt-X-Programm, 2025.

% Weitere Einträge aus REFERENCES.md zu ergänzen.

\end{thebibliography}

\end{document}
